\section{The ASK interface}
\label{sec:interface}
This section describes the decisions, the scorecard rule used in the LongMemEval run, the decision record, and the reference protocol. For each part we state which code produced the reported results and which code is a later specification.

\subsection{Decisions}
Given a stored state and a proposed write (a candidate), ASK returns one of three decisions.
\begin{description}
\item[\keep] The compared prior item, if any, remains active and is not replaced. Whether the candidate is also admitted as an additional active record is a property of the store and not of the decision. The reference protocol's single-slot mode refuses admission, its multi-record mode admits the candidate only when a flag requests it, and the evaluated LongMemEval runtime admits every non-duplicate candidate on \keep{} (11\,464 of 11\,478 \keep{} rows carry a resulting identifier). \keep{} therefore denotes a non-superseding decision and not a withheld write, and a count of \keep{} decisions is not a count of writes that were withheld. In MGP-6, where no store is materialized, \keep{} is a label for the expected preservation of the standing value of the target slot; it tests the decision not to supersede and does not test candidate admission.
\item[\supersede] A specific prior item is replaced by a successor. The predecessor is retained, marked as superseded, and linked to the successor by an identifier and a relation row, within one database transaction in the released code (the memory module used in the evaluated run is not among the files whose hashes the run manifest records).
\item[\abstain] No state change is made. The decision is logged with the candidate and the features that led to it, and the record carries no resulting item identifier.
\end{description}
In the MGP-6 diagnostic the three actions are labels on controlled items and no memory state is materialized (Section~\ref{sec:mgp6}); the two operational meanings of \keep{} above are therefore related by the shared property that no supersession occurs, not by identical state transitions.

\subsection{Scorecard rule}
The controller evaluated on LongMemEval (Section~\ref{sec:lme}) computes five features for a candidate and, when one exists, a selected prior: extraction confidence $c$, evidence quality $e$ (evidence span specificity and a recency comparison), a marker burden $f$ (0.5 for each negation or hedge marker), a signed relation signal $r\in[-1,1]$ from a two-stage language-model relation classifier with deterministic routing (negative values indicate that the candidate contradicts the prior), and a type-conditional stability prior $\rho$ (stable preferences above episodic facts above transient preferences). Two scores are computed from these features, $S_\mathrm{rev}$ with configured weights and $S_\mathrm{abs}$ without weights, and each memory type has thresholds $\theta_\mathrm{sup}$ and $\theta_\mathrm{abs}$. The action rule is
\[
a=\begin{cases}
\abstain & S_\mathrm{abs}>\theta_\mathrm{abs}\\
\supersede & S_\mathrm{rev}>\theta_\mathrm{sup}\ \text{and}\ \max(0,-r)>0\\
\keep & \text{otherwise.}
\end{cases}
\]
Candidates are produced by a language-model extractor, one turn at a time, and the compared prior is selected among the three nearest same-type memories of other sessions by embedding similarity (sentence-transformers/all-MiniLM-L6-v2) using the relation label. The extractor, relation classifier, answer model and judge are language models; the rule above is deterministic given its inputs.

\subsection{Decision record}
Each decision is appended to a JSON-lines action log after the database mutation, with a flush to disk for every row in the released code. A row holds a timestamp, the action, the candidate text and type, a candidate identifier, the session identifier, the prior and resulting memory identifiers, the five features, $S_\mathrm{rev}$, $S_\mathrm{abs}$, both thresholds, and a rationale string. The log is not written atomically with the mutation, so a crash between the two can leave them inconsistent. The record does not contain a configuration digest, an evidence digest or before and after state digests; these fields exist only in the reference protocol below.

\subsection{Reference protocol}
The reference protocol is a Python specification with an abstract store interface and one in-memory implementation. It defines store modes (single-slot versioned, multi-record, and a temporal-event mode that is an enumeration value without an implementation), a transition function that applies a proposal to a store, a transition event that carries evidence, configuration and before and after digests, a decision trace, an invariant validator, declared replay scopes, and two adapters (for the LongMemEval controller and for the MGP-6 live controller). We tested it as follows and report the result without interpretation beyond what the tests show.
\begin{itemize}
\item \textbf{Chronology.} The protocol was first committed on 2026-08-03, after the LongMemEval run (2026-07-18 to 2026-07-19) and after the MGP-6 v4.3 evaluation code (2026-07-25). No module outside the test suite imports it, so no reported number was produced through it.
\item \textbf{Validator.} The five test files that cover the protocol and the memory operations (34 tests) pass at the baseline commit. A full-suite run on a git-archive export of the baseline commit gave 672 passed, 19 skipped and 9 failed; the failures were not investigated one by one, are unrelated to the protocol files in the sampled messages, and are identical with the patch. In an adversarial exercise, the baseline validator accepted all 19 self-consistent tampering probes we constructed, including an altered value or memory type on a retained predecessor, deletion or mutation of an unrelated record, and an unannounced extra record. Separately, the configuration identity is checked for format only: an identity consisting of 64 zeros was accepted when the object was constructed. We then wrote a corrected validator as a patch against the baseline commit (supplement). The same patch was subsequently committed as \texttt{4fc32ca} on a separate local audit branch descending from the baseline commit; it is not part of the baseline commit or of the evaluated runtime. It enforces that a transition changes only the records the action allows (the compared predecessor may only be deactivated, the successor equals the candidate, no unrelated record is removed, changed or added), recomputes event identifiers and the trace evidence digest, and verifies a configuration identity against a canonical SHA-256 of the frozen configuration when the caller supplies that configuration (the constructor itself still checks format only). It rejects all 19 probes and the 64-zero identity with named reasons, accepts eight positive controls, leaves the 34 existing tests passing and adds 98 regression tests (132 passing in the five-file subset plus the new file). The accounting requirement that unrelated records stay unchanged is a design choice of the patch, not a condition that the baseline validator checked, and a store that compacts records would violate it. Limits of the patch: the event identifier is an unkeyed digest, so a self-consistent forgery passes; the supplied before-state is checked only against the event digest; the candidate is supplied by the caller; the configuration digest covers the declared mapping and not the code; the temporal mode has no temporal semantics; and neither the evaluated runtime nor any other store has been run through it.
\item \textbf{Events.} Transition events are returned in memory and are not persisted by the protocol. The evaluated runtime keeps its store in an in-memory SQLite database created for each question, so no store survives a run.
\item \textbf{Scope.} The temporal-event mode is a placeholder. The adapters are reference code and were not used to produce any reported result.
\end{itemize}
We therefore describe the protocol as a specification of the intended record and its checks, with a corrected validator that is a regression-tested reference and not a conformance result, and we claim neither conformance of the evaluated systems nor storage or controller independence as demonstrated properties.
